\documentclass[11pt,letterpaper]{article}
\usepackage[margin=.72in]{geometry}
\usepackage{amsmath,amssymb,amsthm,booktabs}
\usepackage[hidelinks]{hyperref}
\newtheorem{theorem}{Theorem}
\begin{document}
\begin{center}
{\Large\bfseries An exact urn-size estimate from the first repeated draw}\par
Pioneer project \quad 28 November 2026
\end{center}

\begin{abstract}
With uniform sampling from an urn, suppose the first repeat is draw $n+1$.
We prove that the unique maximum-likelihood urn size for $n\geq2$ is
$\binom{n+1}{2}-1-\lfloor(n-1)/3\rfloor$.
A monotone comparison of adjacent likelihoods and two finite binomial
bounds locate the maximum. This proves the empirical formula, second
difference conjecture and generating function recorded in OEIS A111097.
\end{abstract}

Let $a_n$ be the maximum-likelihood urn size when $n$ consecutive draws
with replacement give distinct marbles and the next draw repeats one of
them. The sequence was contributed by Marc A. Brodie. Robert P. P. McKone
conjectured the second differences in 2023. Muslim Alaa recorded the
empirical formula and Stefano Spezia conjectured the generating function
in 2026 \cite{oeis}.

For one-by-one capture and recapture, Leite, Oishi and Pereira proved
uniqueness and characterized the estimate by an adjacent-likelihood
threshold \cite{leite}. Observing $n$ distinct individuals in $n+1$ draws
gives the same likelihood up to a factor independent of the urn size.
Rathie and Z\"ornig gave the ratio and unimodality for the first repeat,
found $a_{100}=5016$, and gave a leading-order approximation
\cite[Example~2.3, equations~(2.7)--(2.9)]{rathie}.
Orlitsky, Santhanam and Viswanathan obtained
$a_n=\binom{n+1}{2}-(n+1)/3+O(1)$ \cite[Lemma~7]{orlitsky}.
We evaluate the threshold explicitly for every $n$.

\begin{theorem}\label{th:urn}
The maximizing urn is unique. We have $a_1=1$, and for $n\geq2$,
\[
 a_n=\binom{n+1}{2}-1-\left\lfloor\frac{n-1}{3}\right\rfloor.
\]
The second differences $a_{n+2}-2a_{n+1}+a_n$ are $2$ at $n=1$ and
$2,1,0$ at residues $0,1,2\pmod3$ for $n\geq2$. Moreover,
\[
 \sum_{n\geq1}a_n z^n
  =\frac{z(1+2z^2-z^3+2z^4-z^5)}{(1-z)^3(1+z+z^2)}.
\]
\end{theorem}

\section{Comparing neighboring urn sizes}
With $N$ marbles the likelihood is zero if $N<n$. Otherwise there are
$N(N-1)\cdots(N-n+1)$ ordered choices for the initial distinct draws and
$n$ choices for the repeat, giving
\[
 L_n(N)=\frac{n\,N(N-1)\cdots(N-n+1)}{N^{n+1}}.
\]
In particular $L_1(N)=1/N$, so $a_1=1$. For $N\geq n$, put $M=N+1$ and
$r=1-1/M$. The ratio $L_n(N+1)/L_n(N)$ is $r^{n+1}/(1-n/M)$.
The sign of the likelihood difference is the sign of
\[
 E(M)=r^{n+1}-(1-n/M)=\frac1M\left(n-\sum_{j=0}^{n}r^j\right).
\]
As $M$ increases, $r$ increases in $[0,1)$, so the expression in parentheses
decreases. It suffices to find the two neighboring integers where the sign
changes.

\section{Locating the sign change}
Write $A=\binom{n+1}{2}$, $s=(n-1)/3$, $t=\lfloor s\rfloor$, and $N_0=A-t-1$.
For $n\geq7$ one has $N_0>n$. For $0\leq x\leq1$, the cubic and quartic
truncations of $(1-x)^{n+1}$ are lower and upper bounds respectively.
Indeed, for $P_k(m,x)=\sum_{j=0}^{k}\binom mj(-x)^j$, Pascal's identity gives
\[
 (1-x)P_k(m,x)=P_k(m+1,x)+(-1)^{k+1}\binom mk x^{k+1}.
\]
Induction in $m$ gives the two bounds by taking $k=3$ and $k=4$.
Substituting $x=1/M$ and using $\binom{n+1}{3}=As$ gives
\begin{equation}\label{eq:bounds}
 \frac{-M^2+AM-As}{M^3}\leq E(M)\leq
 \frac{-M^3+AM^2-AsM+As(n-2)/4}{M^4}.
\end{equation}

At $M=N_0$, the lower numerator in \eqref{eq:bounds} is
$A(t+1-s)-(t+1)^2$. Because $s-t\in\{0,1/3,2/3\}$, it is at least
$A/3-((n+2)/3)^2$. The latter is $(n^2-5n-8)/18>0$ for $n\geq7$.
Thus the likelihood increases from $N_0-1$ to $N_0$.

At $M=N_0+1=A-t$, the upper numerator in \eqref{eq:bounds} is
\[
 -(s-t)AM-t^2M+As(n-2)/4.
\]
If $s=t$, it is negative because
\[
 sM-A(n-2)/4=A(n+2)/12-s^2>0.
\]
Here $A(n+2)/12\geq A/3>((n+2)/3)^2\geq s^2$. If $s-t\geq1/3$, the
numerator is at most $-AM/3+As(n-2)/4<0$, since
\[
 M/3-s(n-2)/4\geq (A-s)/3-s(n-2)/4=(3n^2+11n-2)/36>0.
\]
Thus the likelihood decreases from $N_0$ to $N_0+1$.

For $2\leq n\leq6$, use the integer
$Q_n(N)=N^{n+1}-(N+1-n)(N+1)^n$, which has the same sign as $E(N+1)$.
The comparison below $N_0$ is unnecessary at $n=2$ because $L_2(1)=0$.
\[
\begin{array}{rrrr}
\toprule
n&N_0&Q_n(N_0-1)&Q_n(N_0)\\
\midrule
2&2&\text{--}&-1\\
3&5&6&-23\\
4&8&423&-37\\
5&13&15640&-13607\\
6&19&623579&-2128261\\
\bottomrule
\end{array}
\]
The monotone sign comparison now proves the unique maximum for all $n\geq2$.

\section{Differences and generating function}
For $n\geq2$ the first difference is
$a_{n+1}-a_n=n+1-\mathbf 1_{3\mid n}$. Subtracting consecutive first
differences gives the stated periodic values. At $n=1$, the first three
terms $1,2,5$ give second difference $2$.
Extend by $a_0=0$. If $F(z)=\sum_{n\geq0}a_nz^n$, the coefficients of
$(1-z)^2F(z)$ begin $0,1,0,2$, then repeat the block $0,2,1$.
Multiplying by $1-z^3$ yields $z+2z^3-z^4+2z^5-z^6$, proving the stated
rational generating function.

\begin{thebibliography}{9}
\bibitem{oeis} The OEIS Foundation Inc., \emph{The On-Line Encyclopedia of Integer Sequences},
\href{https://oeis.org/A111097}{entry A111097}.
\bibitem{leite} J. G. Leite, J. Oishi and C. A. B. Pereira,
\emph{Exact maximum likelihood estimate of a finite population size:
capture/recapture sequential sample data}, Probab. Eng. Inf. Sci. 1 (1987),
225--236. \href{https://doi.org/10.1017/S0269964800000425}{doi:10.1017/S0269964800000425}.
\bibitem{rathie} P. N. Rathie and P. Z\"ornig,
\emph{On the birthday problem: some generalizations and applications},
Int. J. Math. Math. Sci. 2003, no. 60, 3827--3840.
\href{https://doi.org/10.1155/S0161171203110101}{doi:10.1155/S0161171203110101}.
\bibitem{orlitsky} A. Orlitsky, N. P. Santhanam and K. Viswanathan,
\emph{Population estimation with performance guarantees},
Proc. IEEE Int. Symp. Inform. Theory (2007), 2026--2030.
\href{https://doi.org/10.1109/ISIT.2007.4557157}{doi:10.1109/ISIT.2007.4557157}.
\end{thebibliography}
\end{document}
